\section{Simulation Studies}
\label{sec:simulation}

This section describes the simulation designs used to generate the results in
\Cref{sec:results}.  The simulations are organized into three groups: foundation
recovery experiments, perturbation-design experiments, and external benchmark
and limitation experiments.

\subsection{Setup}
\label{sec:sim:setup}

In all experiments, data are generated from the full dynamical system, but
inference uses only the post-perturbation equilibrium pairs
$(\tilde{\xv}^{(k)},\uv^{(k)})$.
The main benchmark is the additive polynomial system
\begin{equation}
  \frac{dx_j}{dt}
  =
  \mu_j - \gamma_j x_j
  + \sum_{i\neq j}\{A_{ji}x_i + B_{ji}x_i^2\}
  + u_j,
  \label{eq:sim-additive}
\end{equation}
whose off-diagonal support defines the true directed network; $\gamma_j$ is set
large enough for diagonally stable, positive steady states, and quadratic terms
$B_{ji}$ are activated on a subset of edges in the nonlinear variants.  To
separate transient-model from steady-state-model specification we also use a
multiplicative generalized Lotka--Volterra (gLV) system
$\dot x_j = x_j(r_j + \sum_{i\neq j} A_{ji}x_i - \gamma_j x_j + u_j)$, whose
positive equilibrium reduces to the same additive linear steady-state equation
(\Cref{sec:method:pss}), so PSS-Net is correctly specified for the equilibrium
relation despite the multiplicative transient.

Unless stated otherwise, PSS-Net uses the node-wise \adsiht{} estimator
(\Cref{sec:method:adsiht}) with the no-intercept monomial library through second
order, i.e.\ source-node groups $\{x_i,x_i^2\}$.  Support recovery is summarized
by precision, recall, $F_1$, and Matthews correlation coefficient (MCC) on
off-diagonal directed edges, threshold-free ranking by AUPRC and AUROC, and
effect-strength recovery by coefficient/Jacobian error, sign accuracy, and
edge-function NRMSE.  Summary statistics are reported as mean $\pm$ standard
deviation across 30 independent simulation replicates.

\subsection{Foundation recovery simulations}
\label{sec:sim:foundation}

The foundation simulations use controlled systems in which the true support,
local Jacobian signs, and edge functions are known.  We first simulate additive
and gLV trajectories under a small set of perturbations and retain only their
terminal equilibria as PSS observations.  We then perform an identifiability
scan in which linear, gLV, and additive nonlinear steady-state mechanisms are
fit with linear and quadratic libraries across signal-to-noise ratios.

For the main estimator comparison, node-wise \adsiht{} and group lasso are fit
to shared PSS data under linear and nonlinear truth.  The problem sizes are
$p=8,30,100$, with in-degree $s=2,3,3$ respectively, and the sample budget is
varied through $N/(s\log p)$.  Representative nonlinear fits are used to display
effect decomposition, edge-function recovery, and the inferred directed signed
network.  A separate basis-misspecification experiment crosses quadratic,
saturating Monod-type, and sinusoidal truths with linear, polynomial,
Monod-type, and Fourier-type fitted dictionaries.

\subsection{Scaling and perturbation-design simulations}
\label{sec:sim:design}

The perturbation-design simulations vary both the number and placement of input
conditions.  Sample complexity is summarized by the rescaled budget
$N/(s\log p)$ across multiple network sizes and sparsity levels.  Design
comparisons include random sampling, maximin space-filling designs, oracle
D-optimal designs, and pilot-estimated D-optimal designs.  The D-optimal designs
map candidate perturbations through the known or locally linearized
steady-state response and select an exact subset with
\texttt{AlgDesign::optFederov()}
\citep{wheeler2025algdesign,fedorov1972theory}; the D-optimal backend is a
standard experimental-design tool, not a new inference algorithm.

For the two-stage design experiment, all strategies share the same initial
pilot perturbations, which count toward the total budget.  The pilot-estimated
D-optimal design estimates a local response map from the pilot PSS observations
and then augments the design using the same D-optimal backend.  The resulting
designs are evaluated by running PSS-Net on the selected perturbation-response
pairs and computing the same support-recovery metrics as above.

\subsection{Benchmark, topology, and measurement-limit simulations}
\label{sec:sim:benchmark}

The external benchmark compares PSS-Net with methods that differ in their input
assumptions and output scales.  The benchmark includes modular response analysis
\citep[MRA;][]{kholodenko2002untangling} via the aiMeRA implementation
\citep{jimenezdominguez2021aimera}, sparse linear regression (lasso and elastic
net through \texttt{glmnet}, \citealp{friedman2010glmnet}; group lasso,
\citealp{lounici2011group}), PySINDy-STLSQ
\citep{brunton2016sindy,kaptanoglu2022pysindy}, and association baselines
(Pearson correlation and partial correlation through \texttt{ppcor},
\citealp{kim2015ppcor}).  For methods with different natural score scales,
row-normalized local-response scores are used in the MRA-style comparison, while
threshold-free AUPRC and AUROC are computed from each method's continuous edge
score.

Topology dependence is evaluated by comparing node-wise and joint multi-task
PSS-Net on scale-free and Erd\H{o}s--R\'enyi graphs.  The reported quantities are
edge MCC, out-degree rank correlation, and hub hit rate.  Measurement-scale
sensitivity is evaluated by fitting PSS-Net to four versions of the same
underlying steady states: absolute abundance, closed relative abundance,
centered log-ratio transformed abundance, and relative abundance multiplied by a
noisy total-abundance measurement.
